Un servicio web es una función que un sistema publica en la red para que otro la use por HTTP, sin importar en qué lenguaje está hecho cada uno. El contrato dice qué se envía y qué se recibe. Hay dos estilos: SOAP, que manda mensajes XML dentro de un sobre \cite{w3c_soap12}, y REST, que trabaja con recursos y casi siempre con JSON \cite{fielding2000}.

Para dar idea de la escala: en Prime Day 2026 la API de DynamoDB, que se usa por HTTPS, llegó a 192 millones de peticiones por segundo y pasó de 59 billones en cuatro días \cite{aws_primeday2026}. Un estudio encontró 500 API REST públicas entre los 4000 sitios más visitados, pero solo el 0,8\,\% cumplía todos los principios REST \cite{neumann2021}. En Ecuador tenemos SOAP muy cerca: los comprobantes electrónicos se envían a los servicios de recepción y autorización del SRI \cite{sri_recepcion,sri_autorizacion}. Y en nuestro propio entorno, los tres proyectos del equipo ya se comunican por REST entre el front y el backend; eso lo mostramos en la segunda parte.
